% BTS CIEL 1 — Réseaux — 05 — TCP et UDP
% Compiler deux fois avec XeLaTeX depuis fichiers/.
\documentclass[aspectratio=169,11pt]{beamer}
\usepackage{fontspec}
\setmainfont{Latin Modern Sans}
\setsansfont{Latin Modern Sans}
\setmonofont{Latin Modern Mono}
\usepackage{polyglossia}
\setmainlanguage{french}
\usepackage{tikz}
\usetikzlibrary{arrows.meta,positioning,calc,shapes.geometric}
\usepackage{listings}
\usepackage{xcolor}
\definecolor{foret}{HTML}{1C322D}
\definecolor{creme}{HTML}{F8F3EE}
\definecolor{or}{HTML}{EBB552}
\definecolor{rose}{HTML}{F1CDBE}
\definecolor{vertclair}{HTML}{DDE9DF}
\definecolor{encre}{HTML}{24372E}
\setbeamercolor{normal text}{fg=creme,bg=foret}
\setbeamercolor{frametitle}{fg=or,bg=foret}
\setbeamercolor{title}{fg=creme}
\setbeamercolor{block title}{fg=foret,bg=or}
\setbeamercolor{block body}{fg=encre,bg=creme}
\setbeamercolor{alerted text}{fg=or}
\setbeamercolor{structure}{fg=or}
\setbeamertemplate{navigation symbols}{}
\setbeamertemplate{footline}{\hfill\color{or}\insertframenumber\hspace{0.5cm}\vspace{0.25cm}}
\setbeamertemplate{frametitle}{\vspace{0.35cm}\insertframetitle\par\vspace{0.15cm}\color{or}\rule{\linewidth}{0.6pt}}
\setlength{\parskip}{5pt}
\lstdefinestyle{webcode}{basicstyle=\ttfamily\small\color{encre},keywordstyle=\color{foret},stringstyle=\color{foret},
  backgroundcolor=\color{creme},frame=single,rulecolor=\color{or},framerule=1pt,
  xleftmargin=5pt,xrightmargin=5pt,columns=fullflexible,keepspaces=true,showstringspaces=false}
\newcommand{\question}[1]{\begin{beamercolorbox}[rounded=true,sep=12pt]{block body}\Large\textbf{#1}\end{beamercolorbox}}
\newcommand{\etiquette}[1]{{\color{or}\small\MakeUppercase{#1}}}
% Les schémas ont une surface claire, également lisible en noir et blanc.
\tikzset{
  boite/.style={draw=encre,fill=creme,text=encre,rounded corners,
    align=center,text width=2.8cm,minimum height=1cm},
  flux/.style={-{Stealth},thick,draw=encre},
  every node/.style={text=encre,font=\small}}
\newenvironment{schema}{%
  \begin{beamercolorbox}[sep=8pt,rounded=true]{block body}
  \centering\begin{tikzpicture}[>=Stealth]
}{\end{tikzpicture}\end{beamercolorbox}}

\title{Les protocoles TCP et UDP}
\subtitle{BTS CIEL 1 · Réseaux · 05 — TCP et UDP}
\author{}\date{}
\begin{document}
\begin{frame}[plain]
\vfill\etiquette{Trouver le bon programme}\par\bigskip
{\Huge\bfseries Les protocoles TCP et UDP}\par\bigskip
{\Large BTS CIEL 1 · Réseaux · 05 — TCP et UDP}\vfill
\end{frame}
\begin{frame}{Une machine, plusieurs destinataires}
\question{IP atteint la machine. Quel programme reçoit les données ?}
\begin{schema}
\begin{scope}[scale=1.1,transform shape]
\node[boite,text width=3.5cm,minimum height=2cm,font=\normalsize] (a) at (0,0) {Adresse de\\l'immeuble};
\node[boite,text width=3.5cm,minimum height=2cm,font=\normalsize] (b) at (5,0) {Numéro de\\l'appartement};
\draw[flux] (a)--(b);
\node[font=\normalsize] at (0,-1.5) {IP : machine};
\node[font=\normalsize] at (5,-1.5) {Port : destinataire};
\end{scope}
\end{schema}
\end{frame}
\begin{frame}{Deux connexions vers le même serveur}
\begin{schema}
\node[boite,text width=3.5cm,minimum height=1.5cm,font=\normalsize] (a) at (0,0) {Poste\\\texttt{192.168.10.20}};
\node[boite,text width=3.5cm,minimum height=1.5cm,font=\normalsize] (b) at (8,0) {Serveur\\\texttt{203.0.113.80}};
\draw[flux] (1.9,0.5)--(6.1,0.5);
\node[font=\small] at (4,1.5) {1 : source 50532 → destination 443};
\draw[flux] (1.9,-0.5)--(6.1,-0.5);
\node[font=\small] at (4,-1.2) {2 : source 50533 → destination 443};
\node[font=\normalsize] at (4,-1.8) {Deux connexions TCP distinctes};
\end{schema}
Deux onglets utilisent ici chacun une connexion et un port source.
\end{frame}
\begin{frame}{Lire les ports}
\begin{schema}
\begin{scope}[scale=1.1,transform shape]
\node[boite,text width=3.5cm,minimum height=1.7cm,font=\normalsize] (a) at (0,0) {Client\\port 50532};
\node[boite,text width=3.5cm,minimum height=1.7cm,font=\normalsize] (b) at (7,0) {Serveur\\port HTTPS 443};
\draw[flux] (a)--node[above,fill=creme,font=\small]{1 : demande} (b);
\draw[flux,dashed] (b.south west) to[bend left=20] node[below,fill=creme,font=\small]{2 : réponse vers 50532} (a.south east);
\end{scope}
\end{schema}
\begin{block}{Méthode}
1. Repérez le port source.\\
2. Repérez le port destination.\\
3. Pour la réponse, inversez-les.
\end{block}
\end{frame}
\begin{frame}{Quelques ports usuels}
\begin{beamercolorbox}[sep=10pt,rounded=true]{block body}
\centering\renewcommand{\arraystretch}{1.25}
\begin{tabular}{lll}
\textbf{Service} & \textbf{Transport usuel} & \textbf{Port}\\
DNS & UDP ou TCP & 53\\
DHCP & UDP & 67 serveur, 68 client\\
HTTP / HTTPS & TCP & 80 / 443\\
FTP & TCP & 21 commandes, 20 données\\
SSH & TCP & 22\\
\end{tabular}
\end{beamercolorbox}
Ports bien connus : 0 à 1023. Un port éphémère est un numéro choisi par le client pour une connexion.\\
Sa plage dépend du système.
\end{frame}
\begin{frame}{UDP : envoyer sans connexion}
\question{Faut-il attendre une confirmation pour chaque image vidéo ?}
\begin{schema}
\node[boite,text width=3.3cm,minimum height=2cm,font=\normalsize] (a) at (0,0) {Émetteur\\appel vidéo};
\node[boite,text width=3.3cm,minimum height=2cm,font=\normalsize] (b) at (7,0) {Destinataire\\appel vidéo};
\draw[flux] (a)--node[above,fill=creme,font=\small]{1 : message UDP} (b);
\node[font=\normalsize] at (3.5,-1.7) {Comme une lettre simple : livraison non garantie.};
\end{schema}
UDP ne fournit ni connexion ni accusé de réception. DNS utilise souvent UDP.
\end{frame}
\begin{frame}{TCP : vérifier la livraison}
\question{Que se passe-t-il si 100 octets sont perdus ?}
\begin{schema}
\begin{scope}[scale=0.9,transform shape]
\node[boite,text width=2.5cm,font=\normalsize] at (0,1.3) {Expéditeur};
\node[boite,text width=2.5cm,font=\normalsize] at (8,1.3) {Destinataire};
\draw[thick] (0,0.75)--(0,-2.25); \draw[thick] (8,0.75)--(8,-2.25);
\draw[flux] (0,0.2)--(4,0.2);
\node[above,fill=creme] at (2,0.2) {1 : 100 octets};
\node[font=\normalsize] at (4.8,0.2) {perdu $\times$};
\node at (4,-0.55) {Délai écoulé sans accusé};
\draw[flux] (0,-1.3)--node[above,fill=creme] {2 : retransmission} (8,-1.3);
\draw[flux,dashed] (8,-2.2)--node[above,fill=creme] {3 : accusé de réception} (0,-2.2);
\end{scope}
\end{schema}
TCP remet les octets dans l'ordre.
\end{frame}
\begin{frame}{TCP : ouvrir la connexion}
\begin{schema}
\node[boite,text width=2.7cm,font=\normalsize] at (0,2) {Client};
\node[boite,text width=2.7cm,font=\normalsize] at (8,2) {Serveur};
\draw[thick] (0,1.45)--(0,-2.1); \draw[thick] (8,1.45)--(8,-2.1);
\draw[flux] (0,0.8)--node[above,fill=creme,inner sep=2pt,font=\small]{1 : SYN : demande de connexion} (8,0.8);
\draw[flux,dashed] (8,-0.3)--node[above=2pt,fill=creme,inner sep=2pt,font=\small]{2 : SYN-ACK : accord du serveur} (0,-0.3);
\draw[flux] (0,-1.4)--node[above,fill=creme,inner sep=2pt,font=\small]{3 : ACK : confirmation du client} (8,-1.4);
\end{schema}
Après ces trois messages, les applications peuvent échanger des données.
\end{frame}
\begin{frame}{Choisir TCP ou UDP}
\question{Délai court ou livraison complète : que privilégiez-vous ?}
\begin{schema}
\node[boite,text width=5cm] (q1) at (0,1.5) {Une donnée en retard\\est-elle inutile ?};
\node[boite,text width=5cm] (q2) at (0,-0.5) {Tout doit-il arriver complet\\et dans l'ordre ?};
\node[boite,text width=4.2cm] (udp) at (6.5,1.5) {UDP :\\appel vidéo, audio en direct};
\node[boite,text width=3.2cm] (tcp) at (6.5,-0.5) {TCP : page Web, fichier};
\draw[flux] (q1)--node[above,fill=creme] {oui} (udp);
\draw[flux] (q1)--node[right,fill=creme] {non} (q2);
\draw[flux] (q2)--node[above,fill=creme] {oui} (tcp);
\node at (1,-1.7) {Sinon, étudiez les besoins de l'application.};
\end{schema}
\end{frame}
\begin{frame}{À vous : retrouver une réponse}
\question{À quel port du poste la réponse doit-elle arriver ?}
\begin{schema}
\node[boite,text width=3.5cm,minimum height=2cm,font=\normalsize] (a) at (0,0) {Poste\\\texttt{192.168.10.25}};
\node[boite,text width=3.5cm,minimum height=2cm,font=\normalsize] (b) at (7,0) {Serveur Web\\\texttt{198.51.100.8}};
\draw[flux] (a)--(b);
\node[font=\small] at (3.5,1.5) {1 : TCP source 51000, destination 80};
\node[font=\normalsize] at (3.5,-1.7) {2 : réponse du serveur vers le poste ?};
\end{schema}
\end{frame}
\begin{frame}{Correction : retrouver une réponse}
\begin{schema}
\node[boite,text width=3.5cm,minimum height=2cm,font=\normalsize] (a) at (0,0) {Poste\\\texttt{192.168.10.25}};
\node[boite,text width=3.5cm,minimum height=2cm,font=\normalsize] (b) at (7,0) {Serveur Web\\\texttt{198.51.100.8}};
\draw[flux,dashed] (b)--(a);
\node[font=\small] at (3.5,1.5) {2 : TCP source 80, destination 51000};
\node[font=\normalsize] at (3.5,-1.7) {La réponse retrouve le port éphémère du poste.};
\end{schema}
\end{frame}
\begin{frame}{À vous : choisir un transport}
\question{TCP ou UDP pour ce jeu en ligne ? Pourquoi ?}
\begin{schema}
\node[boite,text width=3cm,minimum height=1.8cm,font=\normalsize] (a) at (0,0) {Joueur};
\node[boite,text width=3cm,minimum height=1.8cm,font=\normalsize] (b) at (7,0) {Serveur de jeu};
\draw[flux] (a)--node[above,fill=creme] {1 : position} (b);
\node[font=\normalsize] at (3.5,-1.7) {30 positions par seconde ; une toutes les 33 ms.};
\end{schema}
Une position perdue est remplacée 33 ms plus tard.
\end{frame}
\begin{frame}{Correction : choisir un transport}
\begin{schema}
\node[boite,text width=3cm,minimum height=1.8cm,font=\normalsize] (a) at (0,0) {Joueur};
\node[boite,text width=3cm,minimum height=1.8cm,font=\normalsize] (b) at (7,0) {Serveur de jeu};
\draw[flux] (a)--node[above,fill=creme] {1 : position UDP} (b);
\node[font=\normalsize] at (3.5,-1.7) {Une nouvelle position arrive 33 ms plus tard.};
\end{schema}
\begin{block}{Choix : UDP}
Une position en retard est inutile. Attendre sa retransmission ralentirait le jeu.
\end{block}
\end{frame}
\begin{frame}{À retenir}
\begin{schema}
\node[boite,text width=3cm,minimum height=1.6cm,font=\normalsize] (a) at (0,0) {IP : la machine};
\node[boite,text width=4cm,minimum height=1.6cm,font=\normalsize] (b) at (6,0) {TCP ou UDP + port :\\le programme};
\draw[flux] (a)--(b);
\end{schema}
\begin{block}{Quatre repères}
IP atteint la machine ; le port désigne le programme.\\
UDP envoie sans connexion et sans garantie de livraison.\\
TCP ouvre une connexion et remet les octets dans l'ordre.\\
Choisissez selon les besoins de l'application.
\end{block}
\end{frame}
\begin{frame}{Sources}
\begin{beamercolorbox}[sep=12pt,rounded=true]{block body}
\textbf{RFC 768} — UDP\\
\texttt{https://www.rfc-editor.org/rfc/rfc768}\\[5pt]
\textbf{RFC 9293} — TCP\\
\texttt{https://www.rfc-editor.org/rfc/rfc9293}\\[5pt]
\textbf{IANA} — registre des ports\\
\texttt{https://www.iana.org/assignments/service-names-port-numbers}\\[5pt]
Support enseignant : « TP DHCP, HTTP, DNS, FTP », p. 1--3.
\end{beamercolorbox}
\end{frame}
\end{document}
